%%% OrchKvCache — SIGMETRICS 2027 Submission
%%% ACM acmsmall format (single-column, ~20 pages)
\documentclass[acmsmall,nonacm,screen]{acmart}

\usepackage{booktabs}
\usepackage{subcaption}
\usepackage{algorithm}
\usepackage{algorithmic}
\usepackage{xcolor}
\usepackage{multirow}
\usepackage{makecell}
\usepackage{balance}

\newcommand{\sys}{OrchKvCache}
\newcommand{\todo}[1]{\textcolor{red}{\textbf{[TODO: #1]}}}
\newcommand{\remark}[1]{\textcolor{blue}{\textbf{[REMARK: #1]}}}

\begin{document}

%% ----------------------------------------------------------------
%%  Title & Authors
%% ----------------------------------------------------------------
\title{\sys{}: Measuring and Scheduling Lossless KV-Cache Residency\\
       Across GPU, DRAM, and SSD}

\author{Ziqing Li}
\affiliation{%
  \institution{Huazhong University of Science and Technology}
  \city{Wuhan}
  \country{China}}
\email{liziqing@hust.edu.cn}

\author{Jianxi Chen}
\affiliation{%
  \institution{Huazhong University of Science and Technology}
  \city{Wuhan}
  \country{China}}
\email{chenjianxi@hust.edu.cn}

%% ----------------------------------------------------------------
%%  Abstract
%% ----------------------------------------------------------------
\begin{abstract}
% Sentence 1 — Problem
Long-context large language models (LLMs) are bottlenecked by KV-cache memory
that grows linearly with sequence length, forcing operators to either over-provision
expensive GPU memory or sacrifice output quality through lossy eviction.
% Sentence 2 — Measurement finding
Through an extensive measurement study across four model families and six
workloads, we discover that attention patterns exhibit \emph{stable hot-set
concentrations} and \emph{predictable reuse distances}, creating an opportunity
for lossless offloading that prior heuristic policies fail to exploit.
% Sentence 3 — System
We present \sys{}, a tiered KV-cache manager that combines reuse-distance
prediction with a critical-path-aware scheduler to place each cache block on the
right tier (GPU, DRAM, or SSD) at the right time, guaranteeing bit-exact
output while hiding promotion latency behind decode-phase compute.
% Sentence 4 — Evaluation result
\todo{Fill final numbers: X$\times$ higher goodput under SLO, Y\% GPU-memory
saving, zero quality degradation across 1000+ prompts on LongBench, RULER,
and AgentBench.}
\end{abstract}

\maketitle

%% ================================================================
%%  1  Introduction
%% ================================================================
\section{Introduction}\label{sec:intro}

\todo{Open with the memory-wall problem for long-context LLM serving: a single
128K-context Llama-3-70B request consumes $>$40\,GB of KV cache, exceeding the
VRAM of a commodity GPU.}

\todo{Introduce the lossless paradox: existing offloading systems either stall
on promotion misses (hurting latency SLO) or fall back to lossy eviction
(hurting quality).  Neither is acceptable for production serving.}

\todo{Preview measurement findings:
\begin{itemize}
  \item Attention hot-sets concentrate on $<$15\% of tokens yet remain stable
        across decode steps.
  \item Block reuse distances follow heavy-tailed distributions that are
        well-predicted by lightweight QK-inner-product proxies.
  \item Promotion latency from DRAM (5--8\,$\mu$s) and SSD (30--80\,$\mu$s)
        can be fully overlapped with the decode-phase compute of most models.
\end{itemize}}

\todo{Describe the design space gap: no prior system combines reuse-distance
modeling, critical-path awareness, and tiered storage in a single framework.}

\todo{List contributions:
\begin{enumerate}
  \item A measurement study characterising attention hot-sets, reuse distances,
        and promotion slack across four model families and six workloads
        (\S\ref{sec:measurement}).
  \item A critical-path residency model that decides \emph{when} and \emph{where}
        to place each KV block with zero quality loss (\S\ref{sec:model}).
  \item The \sys{} system integrating a lightweight signal-acquisition layer,
        a reuse-distance predictor, a critical-path scheduler, and a tiered
        storage backend (\S\ref{sec:design}--\S\ref{sec:impl}).
  \item An evaluation showing X$\times$ goodput improvement under SLO,
        Y\% memory saving, and bit-exact output verified on 1000+ prompts
        (\S\ref{sec:eval}).
\end{enumerate}}


%% ================================================================
%%  2  Background and Motivation
%% ================================================================
\section{Background and Motivation}\label{sec:background}

% --- 2.1 ---
\subsection{The KV Cache Memory Bottleneck}\label{sec:bg:bottleneck}

\todo{Explain KV cache lifecycle in auto-regressive LLM inference: prefill
generates KV for all prompt tokens; decode appends one KV pair per step.
Memory grows as $O(n \cdot d \cdot L)$ where $n$ is sequence length, $d$ is
head dimension, and $L$ is the number of layers.}

\todo{Give concrete numbers: Llama-3-70B at 128K context $\approx$ 40\,GB;
Qwen-2.5-1M at 1M context $\approx$ 300\,GB. Compare with A100-80GB,
H100-80GB VRAM budgets.}

\todo{Discuss operational implications: batch size is memory-capped $\to$
under-utilises compute $\to$ low goodput.  Operators face a cost--quality
trade-off.}

% --- 2.2 ---
\subsection{The Lossless Offloading Paradox}\label{sec:bg:paradox}

\todo{Define lossless: bit-exact reproduction of the attention output that
full-GPU serving would produce.  Any token eviction, quantisation, or
approximation is lossy.}

\todo{Explain why offloading is hard: decode attention needs the \emph{full}
KV set of the attended positions.  If a block is remote when needed, the
decode step stalls until promotion completes.  Stall $\to$ TTFT/TPOT SLO
violation.}

\todo{Survey existing approaches:
\begin{itemize}
  \item \textbf{Lossy eviction}: H2O, StreamingLLM, ScissorHands, SnapKV —
        drop tokens permanently, quality degrades on recall-heavy tasks.
  \item \textbf{Offload-then-prefetch}: InfiniGen, FlexGen, LMCache —
        move cold blocks to DRAM/SSD but rely on coarse heuristics (FIFO,
        LRU) for promotion, causing frequent stalls.
  \item \textbf{Quantised caching}: KIVI, KVQuant — compress KV to fewer
        bits; still lossy, accumulates error over long contexts.
\end{itemize}}

\todo{State the gap: no system provides a principled, measurement-grounded
policy that achieves lossless offloading \emph{without} stalling.}

% --- 2.3 ---
\subsection{Design Space of KV-Cache Management}\label{sec:bg:taxonomy}

\todo{Present a big taxonomy table (Table~\ref{tab:taxonomy}) comparing
axes: eviction vs.\ offloading, lossy vs.\ lossless, policy (FIFO / LRU /
LFU / attention-aware / reuse-distance), tiers (GPU only / GPU+DRAM /
GPU+DRAM+SSD), prefetch strategy (none / speculative / model-based),
critical-path awareness (yes/no).}

\begin{table*}[t]
\centering
\caption{Design space of KV-cache memory management.
\todo{Fill all rows with related-work comparison.}}
\label{tab:taxonomy}
\begin{tabular}{lcccccc}
\toprule
\textbf{System} & \textbf{Evict/Offload} & \textbf{Lossy?} &
  \textbf{Policy} & \textbf{Tiers} & \textbf{Prefetch} &
  \textbf{CP-Aware} \\
\midrule
vLLM~\cite{kwon2023vllm}           & Evict   & Lossless & Seq-level & GPU          & —       & No  \\
H2O~\cite{zhang2024h2o}            & Evict   & Lossy    & Attn-score & GPU         & —       & No  \\
StreamingLLM~\cite{xiao2024streamingllm} & Evict & Lossy & Sink+Window & GPU       & —       & No  \\
InfiniGen~\cite{lee2024infinigen}   & Offload & Lossless & LRU-ish  & GPU+DRAM     & Spec.   & No  \\
FlexGen~\cite{sheng2023flexgen}     & Offload & Lossless & Static   & GPU+DRAM+SSD & None    & No  \\
LMCache~\cite{yao2024lmcache}      & Offload & Lossless & Prefix   & GPU+DRAM+SSD & Prefix  & No  \\
FlexiCache~\cite{li2024flexicache} & Offload & Lossless & Attn-aware & GPU+DRAM   & Pred.   & No  \\
\midrule
\sys{} (ours) & Offload & Lossless & Reuse-dist & GPU+DRAM+SSD & Model & \textbf{Yes} \\
\bottomrule
\end{tabular}
\end{table*}

\todo{Discuss what critical-path awareness means: the scheduler must know
whether a block's next access falls on the critical path of the decode step
and whether promotion latency can be hidden behind compute.}


%% ================================================================
%%  3  Measurement Study
%% ================================================================
\section{Measurement Study}\label{sec:measurement}

\todo{Describe experimental setup for the measurement study: models
(Llama-3-8B, Llama-3-70B, Qwen-2.5-7B, Qwen-2.5-1M-7B), workloads
(LongBench, RULER, AgentBench, needle-in-a-haystack, multi-turn chat,
code generation), instrumented vLLM fork, metrics collected.}

% --- 3.1 ---
\subsection{Attention Concentration and Hot-Set Stability}\label{sec:meas:hotset}

\todo{Key finding: across all models and workloads, $>$85\% of attention mass
concentrates on $<$15\% of KV blocks at each decode step.}

\todo{Show hot-set overlap between consecutive decode steps (Jaccard
similarity $>$0.9 for top-15\% set).  Implication: hot-set membership is
highly predictable from the previous step.}

\todo{Figure~\ref{fig:hotset}: heatmap of per-layer attention concentration
ratio across decode steps for Llama-3-70B on a RULER task.}

\todo{Discuss variation across layers: shallow layers are more diffuse;
deep layers exhibit sharper concentration.  Layer-wise budgeting is
important.}

% --- 3.2 ---
\subsection{Block Reuse Distance Characterisation}\label{sec:meas:reuse}

\todo{Define block-level reuse distance: number of decode steps between
consecutive accesses to the same KV block (or $\infty$ for never-reused
blocks).}

\todo{Key finding: reuse-distance distributions are heavy-tailed (power-law
body, exponential tail) and remarkably stable across sequences of the same
workload type.}

\todo{Figure~\ref{fig:reuse_dist}: CDF of block reuse distances for four
models $\times$ three workloads.  Highlight the ``cold tail'' that comprises
30--60\% of blocks.}

\todo{Implication: a policy that knows the reuse distance of each block can
offload the cold tail with near-zero risk of stall.}

% --- 3.3 ---
\subsection{Promotion Latency and Overlap Slack}\label{sec:meas:slack}

\todo{Measure promotion latency per tier: GPU $\to$ GPU (0), DRAM $\to$ GPU
(5--8\,$\mu$s per block via pinned DMA), SSD $\to$ GPU (30--80\,$\mu$s via
POSIX read, 15--40\,$\mu$s via GDS).}

\todo{Measure decode-step compute time per model: Llama-3-8B $\approx$
X\,ms, Llama-3-70B $\approx$ Y\,ms per step on A100/H100.  This is the
``overlap slack'' available to hide promotion.}

\todo{Key finding: for all tested models, the decode-step compute exceeds
the DRAM promotion latency by $>$10$\times$ and the SSD/GDS promotion
latency by $>$3$\times$, meaning promotion can be fully overlapped with
compute via pipelining.}

\todo{Figure~\ref{fig:slack}: bar chart comparing decode-step compute time
vs.\ promotion latency for each tier and model.}

% --- 3.4 ---
\subsection{Workload Dependence}\label{sec:meas:workload}

\todo{Show how hot-set size, reuse-distance shape, and overlap slack vary
across workload types.  Retrieval-heavy tasks (RULER needle) have sharper
concentration; summarisation tasks have broader attention.}

\todo{Argue that a single static policy cannot cover all workloads $\to$
need an adaptive, measurement-driven approach.}


%% ================================================================
%%  4  Critical-Path Residency Model
%% ================================================================
\section{Critical-Path Residency Model}\label{sec:model}

\todo{High-level goal: given the measured signals, derive a per-block
placement decision that guarantees (a) lossless output and (b) zero
critical-path stall, while minimising GPU memory usage.}

% --- 4.1 ---
\subsection{Reuse Distance Prediction}\label{sec:model:reuse}

\todo{Formalise the reuse-distance prediction problem.  Input: QK
inner-product proxy score $s_i^{(t)}$ for block $i$ at step $t$,
historical reuse distances.  Output: predicted reuse distance
$\hat{r}_i^{(t)}$.}

\todo{Present the predictor: lightweight exponential moving average of
per-block proxy scores, thresholded to classify blocks into hot / warm /
cold tiers.  Justify via correlation analysis from \S\ref{sec:meas:reuse}.}

\todo{Complexity: $O(B)$ per decode step where $B$ is the number of active
blocks.  No neural network overhead.}

% --- 4.2 ---
\subsection{Promotion Latency Model}\label{sec:model:latency}

\todo{Model the promotion latency $\ell_{\text{tier}}(b)$ for block $b$
from each tier as a function of block size, bus bandwidth, queue depth, and
concurrency.  Calibrate with micro-benchmarks.}

\todo{Define overlap slack $\Delta^{(t)}$: the compute time of the current
decode step available for hiding promotion.  $\Delta$ depends on model size,
batch size, and hardware.}

% --- 4.3 ---
\subsection{Offload Condition and Tier Selection}\label{sec:model:offload}

\todo{Derive the offload condition: block $b$ can be offloaded to tier $T$
iff its predicted reuse distance $\hat{r}_b > \lceil \ell_T(b) / \Delta
\rceil + \text{safety margin}$.  This guarantees that promotion completes
before the block is needed.}

\todo{Tier selection: prefer DRAM for warm blocks (moderate reuse distance),
SSD for cold blocks (long reuse distance).  Tie-break by promotion latency
ratio.}

\todo{Handle prediction errors: if a block is accessed earlier than
predicted, the scheduler has a ``rescue'' path that preempts lower-priority
promotions and uses emergency synchronous DMA.  Quantify rescue rate from
measurement traces.}

% --- 4.4 ---
\subsection{Model Validation}\label{sec:model:validation}

\todo{Validate the model on measurement traces: replay recorded access
patterns through the residency model and compare predicted stall vs.\ actual
stall.  Show $<$2\% prediction error on median stall time.}

\todo{Figure~\ref{fig:model_validation}: scatter plot of predicted vs.\
measured stall per decode step across four workloads.}


%% ================================================================
%%  5  System Design
%% ================================================================
\section{System Design}\label{sec:design}

% --- 5.1 ---
\subsection{Architecture Overview}\label{sec:design:arch}

\todo{Figure~\ref{fig:arch}: system architecture diagram showing four
components: (1) Signal Acquisition Layer, (2) Reuse-Distance Predictor,
(3) Critical-Path Scheduler, (4) Tiered Storage Backend.  Show data flow
from attention kernel $\to$ predictor $\to$ scheduler $\to$ DMA/GDS engine.}

\todo{Explain the per-step control loop: after each decode step, the signal
layer emits proxy scores; the predictor updates reuse estimates; the
scheduler issues offload / promote / pin commands; the backend executes
asynchronous transfers.}

% --- 5.2 ---
\subsection{Low-Overhead Signal Acquisition}\label{sec:design:signal}

\todo{QK proxy: compute partial QK inner products using only the query of the
current decode step and a sampled subset of K heads.  Cost: one small GEMV
per layer, $<$1\% of decode compute.}

\todo{Sampling: only sample every $S$ decode steps (default $S{=}4$) and
interpolate between samples.  Reduces overhead by $4\times$ with $<$0.5\%
prediction accuracy loss.}

\todo{Bypass: for the first $P$ prefill tokens (where attention is dense by
construction), skip signal acquisition entirely.}

% --- 5.3 ---
\subsection{Reuse-Distance-Aware Residency Predictor}\label{sec:design:predictor}

\todo{Describe the EMA-based predictor from \S\ref{sec:model:reuse} in
implementation detail.  Per-block state: 4 bytes (16-bit EMA score + 16-bit
last-access timestamp).  Total overhead for 1M blocks: 4\,MB.}

\todo{Explain the hot / warm / cold classification thresholds and how they
adapt to workload-specific reuse-distance distributions using online
percentile tracking.}

% --- 5.4 ---
\subsection{Critical-Path-Aware Scheduler}\label{sec:design:scheduler}

\todo{Describe the scheduling algorithm:
\begin{enumerate}
  \item Classify blocks into hot (keep on GPU), warm (offload to DRAM),
        cold (offload to SSD).
  \item For each warm/cold block, compute the latest safe offload time
        using the offload condition from \S\ref{sec:model:offload}.
  \item Schedule promotions to start $\hat{r}_b - \lceil \ell_T / \Delta
        \rceil$ steps before predicted access.
  \item If multiple promotions compete for bandwidth, prioritise by
        soonest predicted access.
\end{enumerate}}

\todo{Discuss concurrency management: cap in-flight DMA to avoid
bus saturation; back-pressure mechanism when promotion queue is full.}

% --- 5.5 ---
\subsection{Tiered Storage Backend}\label{sec:design:backend}

\todo{Three I/O paths:
\begin{itemize}
  \item \textbf{DRAM tier}: pinned host memory, \texttt{cudaMemcpyAsync}
        via CUDA streams.
  \item \textbf{SSD tier (POSIX)}: page-cache-bypassing \texttt{O\_DIRECT}
        reads into pinned buffers, then DMA to GPU.
  \item \textbf{SSD tier (GDS)}: NVIDIA GPUDirect Storage for direct
        SSD$\to$GPU transfers, eliminating CPU bounce buffer.
\end{itemize}}

\todo{Block layout on SSD: contiguous 2\,MB blocks aligned to filesystem
block size.  Pre-allocated file pool to avoid allocation jitter.}

\todo{Discuss GDS availability and fallback to POSIX on systems without GDS
support.}

% --- 5.6 ---
\subsection{Production Integration Path}\label{sec:design:production}

\todo{Describe how \sys{} integrates with vLLM's PagedAttention:
replace the block allocator with a tiered allocator; hook into the scheduler
to receive offload/promote commands; minimal kernel modification (only the
signal acquisition GEMV).}

\todo{Discuss deployment considerations: multi-GPU (tensor-parallel),
multi-node, and disaggregated prefill--decode architectures.}


%% ================================================================
%%  6  Implementation
%% ================================================================
\section{Implementation}\label{sec:impl}

\todo{Lines of code: \sys{} is implemented in $\sim$X\,K lines of Python
and $\sim$Y\,K lines of CUDA/C++.}

\todo{Built on top of vLLM v0.6.x.  Key modifications:
\begin{itemize}
  \item Tiered block allocator (Python, $\sim$X lines).
  \item Signal acquisition CUDA kernel (Triton, $\sim$Y lines).
  \item Reuse-distance predictor (Python + NumPy, $\sim$Z lines).
  \item Critical-path scheduler (Python, $\sim$W lines).
  \item GDS backend (C++ extension, $\sim$V lines).
\end{itemize}}

\todo{Discuss engineering challenges: CUDA stream synchronisation for
overlapped promotion, memory pool fragmentation, GDS alignment requirements.}


%% ================================================================
%%  7  Evaluation
%% ================================================================
\section{Evaluation}\label{sec:eval}

% --- 7.1 ---
\subsection{Experimental Setup}\label{sec:eval:setup}

\todo{Hardware: NVIDIA A100-80GB / H100-80GB GPUs, 512\,GB DDR5 DRAM,
Samsung 990 PRO NVMe SSD (7\,GB/s seq read).  Server: 2$\times$ AMD EPYC
9654, 192 cores.}

\todo{Models: Llama-3-8B, Llama-3-70B (TP=4), Qwen-2.5-7B, Qwen-2.5-1M-7B.}

\todo{Workloads: LongBench, RULER, AgentBench, needle-in-a-haystack,
multi-turn chat (ShareGPT), code generation (HumanEval-X with 32K context).}

\todo{Baselines: vLLM (no offload), FlexGen, InfiniGen, LMCache, FlexiCache.
Policy baselines: FIFO, LRU, LFU, Belady (oracle).}

\todo{Metrics: goodput (requests/sec meeting SLO), TTFT, TPOT, GPU memory
usage, output quality (exact-match, F1, BLEU).}

% --- 7.2 ---
\subsection{Model Validation: Predicted vs.\ Measured Stall}\label{sec:eval:validation}

\todo{Figure: scatter plot of predicted vs.\ measured per-step stall time
for \sys{}'s residency model.  Show $R^2 > 0.95$ across workloads.}

\todo{Table: mean absolute error of stall prediction per model and
workload.  Target: $<$5\% MAPE.}

% --- 7.3 ---
\subsection{Policy Comparison}\label{sec:eval:policy}

\todo{Compare seven policies under the same tiered backend:
FIFO, LRU, LFU, Belady (oracle upper bound), InfiniGen-style (speculative
QK), FlexiCache-style (attention-score ranking), \sys{} (reuse-distance +
critical-path).}

\todo{Figure: goodput vs.\ GPU memory budget (10\%--100\% of full KV) for
each policy on Llama-3-70B / RULER.}

\todo{Key result: \sys{} tracks within X\% of Belady and outperforms the
best heuristic (FlexiCache-style) by Y\%.}

% --- 7.4 ---
\subsection{System Comparison}\label{sec:eval:system}

\todo{End-to-end comparison with vLLM, LMCache, FlexGen, InfiniGen under
realistic serving load (Poisson arrivals, mixed context lengths).}

\todo{Figure: goodput vs.\ request rate for each system at 50\% GPU memory
budget.  \sys{} sustains X$\times$ higher goodput than vLLM and Y$\times$
higher than the best offloading baseline.}

\todo{Table: TTFT p50/p99 and TPOT p50/p99 comparison.}

% --- 7.5 ---
\subsection{Goodput Under SLO vs.\ Request Rate}\label{sec:eval:slo}

\todo{Figure: SLO attainment rate (fraction of requests meeting TPOT
$<$100\,ms) vs.\ request rate for each system.}

\todo{Show that \sys{} maintains $>$95\% SLO attainment at X req/s while
baselines drop below 80\% at Y req/s.}

% --- 7.6 ---
\subsection{Memory Pressure Sweep}\label{sec:eval:memory}

\todo{Vary GPU KV-cache budget from 5\% to 75\% of full KV.  Plot goodput
and TPOT for each system.}

\todo{Key finding: \sys{} degrades gracefully; at 10\% budget it still
achieves Z\% of full-GPU goodput.  Baselines cliff at $<$30\% budget.}

% --- 7.7 ---
\subsection{SSD Backend Ablation}\label{sec:eval:ssd}

\todo{Compare three I/O backends: (a) GPU+DRAM only, (b) GPU+DRAM+SSD
(POSIX O\_DIRECT), (c) GPU+DRAM+SSD (GDS).}

\todo{Show that adding SSD tier extends effective KV capacity by X$\times$
with $<$Y\% goodput loss compared to DRAM-only.  GDS further improves SSD
tier throughput by Z\%.}

% --- 7.8 ---
\subsection{Signal Acquisition Overhead}\label{sec:eval:overhead}

\todo{Micro-benchmark: measure the per-step overhead of the QK proxy kernel
and predictor update as a fraction of decode-step compute.}

\todo{Target: $<$1\% overhead on Llama-3-70B, $<$2\% on Llama-3-8B (smaller
model $\to$ smaller denominator).}

\todo{Ablation on sampling interval $S$: show accuracy vs.\ overhead
trade-off.}

% --- 7.9 ---
\subsection{Correctness and Quality}\label{sec:eval:quality}

\todo{Run 1000+ prompts from LongBench, RULER, and AgentBench through
both full-GPU vLLM and \sys{}.  Compare outputs token-by-token.}

\todo{Result: 100\% exact match (bit-exact output), confirming lossless
guarantee.}

\todo{Table: benchmark scores (F1, exact-match, BLEU, pass@1) for
full-GPU vs.\ \sys{} — identical rows.}


%% ================================================================
%%  8  Related Work
%% ================================================================
\section{Related Work}\label{sec:related}

\todo{Organise into subsections or paragraphs:}

\paragraph{KV-Cache Eviction and Compression.}
\todo{H2O~\cite{zhang2024h2o}, StreamingLLM~\cite{xiao2024streamingllm},
ScissorHands~\cite{liu2024scissorhands}, SnapKV~\cite{li2024snapkv},
PyramidKV~\cite{cai2024pyramidkv}, Keyformer~\cite{adnan2024keyformer},
Ada-KV~\cite{feng2024adakv}, KIVI~\cite{liu2024kivi},
PagedEviction~\cite{yu2024pagedeviction}.
These are all lossy — \sys{} is lossless.}

\paragraph{KV-Cache Offloading.}
\todo{FlexGen~\cite{sheng2023flexgen}, InfiniGen~\cite{lee2024infinigen},
LMCache~\cite{yao2024lmcache}, FlexiCache~\cite{li2024flexicache},
CachedAttention~\cite{gao2024cachedattention},
SqueezeAttention~\cite{yang2024squeezeattention},
CacheBlend~\cite{yao2024cacheblend},
vTensor~\cite{xu2024vtensor}, Tokencake~\cite{patel2024tokencake}.
Offload but lack critical-path-aware scheduling.}

\paragraph{GPU--Storage Systems.}
\todo{BaM~\cite{qureshi2023bam}, ALISA~\cite{zhao2024alisa},
KVDrive~\cite{patel2024kvdrive}, GDS/TERAIO~\cite{nvidia2023gds}.
Hardware-level I/O but no KV-cache-specific policy.}

\paragraph{Attention Kernels and Serving Systems.}
\todo{FlashAttention~\cite{dao2022flashattention},
FlashAttention-2~\cite{dao2023flashattention2},
FlashAttention-3~\cite{shah2024flashattention3},
FlashInfer~\cite{ye2024flashinfer},
vLLM~\cite{kwon2023vllm},
Mooncake~\cite{qiu2024mooncake},
Tutti~\cite{wu2024tutti},
Quest~\cite{tang2024quest},
DUAL-BLADE~\cite{wang2024dualblade}.
Complementary to \sys{} — our system can layer on top of any attention kernel.}

\paragraph{KV-Cache Quantisation and Other Compression.}
\todo{KIVI~\cite{liu2024kivi}, TurboQuant~\cite{chen2024turboquant},
CommVQ~\cite{yang2024commvq}, IMPRESS~\cite{liu2024impress}.
Lossy; can be composed with \sys{} as an optional tier.}


%% ================================================================
%%  9  Discussion and Limitations
%% ================================================================
\section{Discussion and Limitations}\label{sec:discussion}

\todo{Discuss limitations:
\begin{itemize}
  \item Prediction accuracy degrades for highly irregular attention patterns
        (e.g., adversarial prompts).
  \item GDS requires specific NVMe hardware and driver support.
  \item Current prototype is single-node; multi-node extension is future work.
  \item Overhead of signal acquisition may become significant for very small
        models ($<$1B parameters).
\end{itemize}}

\todo{Discuss broader impact: enabling longer-context serving on commodity
hardware democratises access to powerful LLMs.}

\todo{Discuss potential negative societal impacts and mitigations (standard
ACM ethics).}


%% ================================================================
%%  10  Conclusion
%% ================================================================
\section{Conclusion}\label{sec:conclusion}

\todo{Restate the problem, the key insight from the measurement study, the
system design, and the main evaluation results in 4--5 sentences.}

\todo{Highlight that \sys{} achieves X$\times$ goodput improvement and Y\%
memory saving while maintaining bit-exact output quality, validated on
1000+ prompts across diverse benchmarks.}

\todo{Mention future directions: multi-node tiered caching, integration with
disaggregated serving architectures, and co-design with attention-kernel
compression.}


%% ================================================================
%%  Acknowledgements
%% ================================================================
\begin{acks}
\todo{Acknowledge funding sources, compute providers, and colleagues who
provided feedback.}
\end{acks}


%% ================================================================
%%  References
%% ================================================================
\bibliographystyle{ACM-Reference-Format}
\bibliography{references}

%% ================================================================
%%  Appendix (optional)
%% ================================================================
\appendix

\section{Extended Measurement Data}\label{app:measurement}
\todo{Additional figures and tables from the measurement study that do not
fit in the main body.}

\section{Proof of Offload Safety Condition}\label{app:proof}
\todo{Formal proof that the offload condition in \S\ref{sec:model:offload}
guarantees zero stall under the stated assumptions.}

\section{Artifact Description}\label{app:artifact}
\todo{Describe the artifact for reproducibility: code repository, Docker
image, scripts to reproduce all figures and tables.}

\end{document}
